\documentclass[11pt]{article}
\usepackage[margin=1in]{geometry}
\usepackage{amsmath,amssymb,booktabs}
\usepackage[hidelinks]{hyperref}

\title{Learnable Stochastic Differential Equations for Movement Prediction}
\author{Author information to be added after approval}
\date{Draft --- no empirical claims included}

\begin{document}
\maketitle

\begin{abstract}
Movement prediction for search and rescue must represent uncertainty rather
than returning a single nominal path. This draft specifies a research paper
around a learnable stochastic differential equation (SDE) implementation and a
separate Lean formalization of selected probability identities. The executable
baseline is a segment-level latent-regime underdamped Langevin model with
exact-Gaussian, Euler--Maruyama, split-step, and common-random-number inference
routes. The formal development checks algebraic statements for soft evidence
and mixture weights. No private trajectory data, trained checkpoint, empirical
metric, or performance conclusion is reported here. The experimental sections
state a protocol to be completed only with approved, versioned aggregate
evidence.
\end{abstract}

\section{Introduction}

Search planning benefits from a forecast distribution over plausible movement
states at one or more future horizons. Such a distribution must retain the
uncertainty induced by incomplete observations and heterogeneous movement
behaviour. A stochastic differential equation provides a direct representation
of state evolution while allowing the forecast to be sampled or summarized by
moments.

This paper documents the present research artefacts without extending their
claims. The software repository implements a typed Python framework for model
fitting and inference. Its current executable baseline uses a latent regime
that remains fixed inside a trajectory segment. The Lean repository formalizes
a finite-dimensional, finite-measure algebraic core for soft evidence. Analytic
existence, uniqueness, and Doob \(h\)-transform results are outside its stated
formalization scope.

The paper will make three bounded contributions once the required evidence is
available: an auditable implementation description, a precise account of the
machine-checked theorem boundary, and a preregistered evaluation protocol. It
does not claim that the current model improves any search-and-rescue outcome.

\section{Problem formulation}

Let \(X_t=(p_t,v_t)^\top\in\mathbb{R}^2\) denote a position--velocity state
for one coordinate of a movement segment. A contextual model may be written as

\begin{equation}
  dX_t=b_\theta(t,X_t,c_t)\,dt+\sigma_\theta(t,X_t,c_t)\,dW_t,
  \label{eq:general-sde}
\end{equation}

where \(c_t\) contains approved contextual inputs and \(W_t\) is a Wiener
process. The released repository does not include a registered neural SDE or a
public context-data pipeline. Equation~\eqref{eq:general-sde} therefore states
the modelling interface, not an empirical model claim.

Given an observed initial state, the forecast task returns samples or moments
of \(X_{t+\tau}\) for requested horizons \(\tau\). Evaluation must use splits
defined outside the source repository and report only approved aggregate
statistics.

\section{Implemented segment-level model}

The current baseline is an affine underdamped Langevin SDE with a latent regime
\(r\in\{1,\ldots,K\}\). The regime is supplied explicitly and is constant
within each segment. Its dynamics are

\begin{align}
  dp_t &= v_t\,dt, \\
  dv_t &= \left[-\Gamma_r v_t+(a_r-\kappa)p_t+c_r\right]dt+g_r\,dW_t.
  \label{eq:segment-sde}
\end{align}

Here \(\Gamma_r\), \(a_r\), \(c_r\), and \(g_r\) are regime-specific
parameters, while \(\kappa\) is a model parameter shared by regimes. The
state transition is affine Gaussian, so the implementation provides an exact
transition kernel when this model is selected.

The estimation code initializes segment labels from features derived from
velocity and displacement variation. It then alternates segment posterior
weights with weighted least-squares updates of a discrete transition, followed
by a discrete-to-continuous conversion. This is an implementation description;
convergence of the estimator on a target dataset remains an empirical question.

\section{Inference and software boundary}

The Python framework separates the SDE model, estimator, inference engine,
evaluation rule, and artifact adapters. Exact-Gaussian sampling is available
for models that provide an exact transition. Euler--Maruyama is an explicitly
selected numerical route. The segment-level model also has a split-step route,
and common random numbers can be used for shared-path, multi-horizon sampling.
Unsupported combinations fail explicitly rather than silently selecting another
method.

The repository's synthetic workflow is deterministic and is intended for
installation and regression tests. It is not evidence for real-data predictive
accuracy. Neural-SDE and Fokker--Planck variants remain unregistered research
skeletons, and condition-feature alignment has not fully moved into the
declarative data-source layer.

\section{Soft evidence and formal verification boundary}

For binary soft evidence with prior \(\pi\), likelihoods \(\ell_1\) and
\(\ell_0\), the formal development includes the posterior mixture weight

\begin{equation}
  w=\frac{\ell_1\pi}{\ell_1\pi+\ell_0(1-\pi)}.
  \label{eq:posterior-weight}
\end{equation}

It checks the mixture identity, a likelihood-ratio form of
Equation~\eqref{eq:posterior-weight}, a total-variation mixture bound, and
properties of the logistic weight. The Lean development also covers a product
identity for evidence-pattern weights and a nested-evidence concentration
result under the assumptions encoded in its theorem signature.

These results are not a formal proof of SDE well-posedness or of a Doob
\(h\)-transform for the target process. The manuscript must preserve this
boundary: only declarations present in the Lean source and its axiom audit may
be described as machine checked.

\section{Data governance and reproducibility}

The paper repository must not contain raw or transformed trajectories,
coordinates, timestamps, identifiers, row-level predictions, non-public
checkpoints, or archives containing such material. The current software
contract expects a locally controlled data root with a main trajectory table
and a split manifest. Its physical authoritative location and data steward are
not established by this paper draft.

Every reported result should record a logical data version, split-manifest
version, source revision of the Python and Lean repositories, random-seed
policy, model configuration, and aggregate-only output. A controlled runner may
mount the approved data root read-only. Public continuous integration should use
only deterministic synthetic data.

\section{Experimental protocol}

This section is intentionally a protocol, not a result section. Before any
table or conclusion is added, the following items must be fixed and reviewed:

\begin{enumerate}
  \item the authorised dataset snapshot, split manifest, licence, and privacy classification;
  \item forecast horizons, observation window, state construction, and missing-data handling;
  \item baseline models and the exact inference configuration for each comparison;
  \item aggregate scoring rules, including one explicit energy-score convention;
  \item the number of seeds, uncertainty intervals, and the rule for handling failed runs;
  \item a disclosure of known numerical limitations and any deviation from the preregistered plan.
\end{enumerate}

Table~\ref{tab:results-template} is a publication template. Its cells must
remain blank until values are generated by an approved run and independently
checked against the recorded configuration.

\begin{table}[ht]
\centering
\caption{Aggregate evaluation template. Values are intentionally absent.}
\label{tab:results-template}
\begin{tabular}{llll}
\toprule
Model & Horizon & Aggregate score & Run reference \\
\midrule
Segment-level SDE & [to be specified] & [not yet available] & [to be specified] \\
Approved baseline & [to be specified] & [not yet available] & [to be specified] \\
\bottomrule
\end{tabular}
\end{table}

\section{Limitations and conclusion}

The current artefacts establish an executable baseline, explicit numerical
routing, deterministic synthetic tests, and a bounded formalization. They do
not establish predictive performance on private trajectories, operational value
for a search mission, or formal analytic properties beyond the Lean scope. The
next revision should add only reviewed aggregate evidence and must update the
limitations when the software or formal boundary changes.

\paragraph{Source-artifact note.} Technical statements in this draft are
traceable to the \href{https://github.com/HuAzurestar/learnable-sde-for-movement-prediction}{Python implementation repository}
and the \href{https://github.com/HuAzurestar/learnable-sde-theory-to-predict-movement}{Lean formalization repository}.
Bibliographic citations have deliberately not been invented; verified references
belong in \texttt{references.bib} before submission.

\end{document}
